% 6.电路基础.tex 子文件，book文档类，主文档通过\input{6.电路基础.tex}导入
\chapter{第三章 电路定理}

\section{置换定理}
置换定理：在任意线性和非线性电路中，若已知某一端口的电压和电流分别为$U$和$I$，则可用$U_S=U$的电压源或$I_S=I$的电流源来置换此一端口，而不影响电路中其他部分的电压和电流。

在应用置换定理时还需说明：
\begin{enumerate}
    \item 置换定理要求置换后的电路有唯一解
    \item 除被置换部分发生变化外，其余部分在置换前后必须保持完全相同
\end{enumerate}
根据置换定理可以得出推论：

\section{齐性定理和叠加定理}
\subsection{齐性定理}
齐性定理：在只有一个激励$X$作用的线性电路中，设任一响应为$Y$，记作$Y=f(X)$，若将该激励乘以常数$K$，则对应的响应$Y'$也等于原来响应乘以同一常数，即
\[
Y'=f(KX)=Kf(X)=KY
\]
从数学上看。

对于线性直流电路，其电路方程为线性代数方程，此时齐性定理可直观表述为：若电路中只有一个激励，则响应与激励成正比。

\subsection{叠加定理}
叠加定理：在线性唯一解电路中，有几个独立电源共同作用产生的响应等于各个独立电源单独作用时产生响应的代数叠加。

值得注意的是，功率与独立电源的源电压和源电流不是线性关系，所以功率不能直接用叠加关系来计算。

对于线性直流电路。

\section{等效电源定理}
\subsection{戴维宁定理}
线性含源一端口网络的对外作用可以用一个电压源串联电阻的电路来等效代替。其中电源的源电压等于此一端口网络的开路电压，而电阻等于此一端口网络内部独立电源置零后所得无独立源一端口网络的等效电阻。

\subsection{诺顿定理}

\section{特勒根定理}
特勒根定理：在电路。

\section{互易定理}
定理（第一种形式）：对于由一个独立电压源和若干线性二端电阻组成的电路，当此电压源。

定理（第二种形式）：

定理（第三种形式）：

\section{对偶原理}
如果。

\endinput